\begin{helper}
Let \(p,k,t\) be natural numbers with \(0<k\) and \(t\in\{1,\ldots,p\}\).
Then the number of bad \(2k\)-tuples over \(\mathbb F_2^p\) whose final
first-coordinate span rank is \(t\) is at most
\[
  \binom{k}{t}\,
  2^{\,p(t+k)-\binom{t+1}{2}}.
\]
Equivalently,
\[
\left|\{ab:F2BadTuple(p,k,ab)\ \text{and}\
  F2BadTupleRank(p,k,ab,k)=t\}\right|
\le
\binom{k}{t}2^{\,p(t+k)-\binom{t+1}{2}}.
\]
\end{helper}
\begin{proof}
For a bad tuple \(ab\) with final rank \(t\), define its Boolean increase
pattern \(z_{ab}:\{0,\ldots,k-1\}\to\{\mathrm{false},\mathrm{true}\}\) by
\[
  z_{ab}(i)=\mathrm{true}
  \quad\Longleftrightarrow\quad
  F2BadTupleRank(p,k,ab,i+1)=F2BadTupleRank(p,k,ab,i)+1.
\]
The rank-increase set of \(\noderef{F2BadTupleRankIncreaseSet}\) is exactly
the set of indices where this Boolean sequence is true. Since \(ab\) is bad
and \(0<k\), \(\noderef{F2BadTupleRankIncreaseSetCard}\) says that this
increase set has cardinality \(F2BadTupleRank(p,k,ab,k)\). For the tuples
being counted here the final rank is \(t\), so
\[
  \operatorname{wt}(z_{ab})=t.
\]

Now send such a tuple \(ab\) to the pair consisting of its increase pattern
\(z_{ab}\) and the same tuple \(ab\), viewed as an element of the fixed-pattern
fiber for \(z_{ab}\). This map is injective because its second component is
the original tuple. Hence the final-rank fiber has cardinality at most the
cardinality of the disjoint union, over all Boolean patterns \(z\) of weight
\(t\), of the fixed-pattern fibers counted by
\(\noderef{F2BadTupleFixedIncreaseCount}\).

Fix one such pattern \(z\). Since \(\operatorname{wt}(z)=t\) and \(t\le p\),
\(\noderef{F2BadTupleFixedIncreaseCountBound}\) gives
\[
  F2BadTupleFixedIncreaseCount(p,k,z)
  \le
  2^{\,p(t+k)-\binom{t+1}{2}}.
\]
Thus every fixed-pattern fiber in the disjoint union has size at most this
same bound. By \(\noderef{BinarySequenceWeightFiberCard}\), the number of
Boolean patterns of weight \(t\) is \(\binom{k}{t}\). Multiplying the common
fiber bound by the number of possible patterns gives
\[
  \binom{k}{t}\,
  2^{\,p(t+k)-\binom{t+1}{2}},
\]
which is the desired estimate.
\end{proof}
